\documentclass[../../../uem_mei_q.tex]{subfiles}

\begin{document}
    $n$を自然数とし，座標平面上で異なる$2n$個の点\\
    $\mathrm{P}_1,\:\mathrm{P}_2,\:\cdots,\:\mathrm{P}_{2n}$をとる．%
    この$2n$個の点の中から2点ずつを結んでいき，次の条件を満たす$n$個の線分をつくる．
    
    \begin{itemize}[leftmargin=5\zw, labelsep=.7\zw, topsep=4pt]
        \item[（条件）] $n$個の線分の端点は，全て異なる．
    \end{itemize}
    ここでは，$n$個の線分を選ぶといったら，常にこの条件を満たすものとする．%
    $n$個の線分を選んだとき，その$n$個の線分の長さの和を，$\ell$とする．%
    また，$\ell$が最小となる$n$個の線分の選び方の総数を，$a_n$と置く．\\
    このとき，次の問に答えよ．

    \begin{myenuc}
        \item 次の命題の内，正しい命題をひとつ選び，そのアルファベットを，解答欄に記入せよ．
        
        \begin{myenua}[label=\Alph*., leftmargin=2\zw]
            \item $\ell$が最小となる$n$個の線分を選んだとき，どの2つの線分も共有点を持たない．
            \item $\ell$が最小となる$n$個の線分の選び方の必要十分条件は，%
                $n$個の線分のどの2つの線分も共有点を持たないことである．
            \item ある2つの線分が，共有点を持つが，$\ell$が最小となる$n$個の線分の選び方がある．
            \item $a_n=1$ ならば，$\mathrm{P}_1,\:\mathrm{P}_2,\:\cdots,\:\mathrm{P}_{2n}$ %
                は，直線上にある．
            \item $\mathrm{P}_1,\:\mathrm{P}_2,\:\cdots,\:\mathrm{P}_{2n}$ が，%
                直線上にある為の必要十分条件は，$a_n=1$ である．
        \end{myenua}
        \item $\mathrm{P_1}(1,\;1),\:\mathrm{P_2}(1,\;2),\:\cdots,\:\mathrm{P_n}(1,\;n),%
            \:\mathrm{P_{n+1}}(2,\;1),\\
            \mathrm{P_{n+2}}(2,\;2),%
            \:\cdots,\:\mathrm{P_{2n}}(2,\;n)$ の場合に，次の問に答えよ．
            \begin{myenud}
                \item $a_2$，$a_3$ を求めよ．（答えのみを，解答欄に記入せよ．）
                \item 一般の$n$の場合に，$\ell$の最小値を求めよ．（答えのみを，解答欄に記入せよ．）
                \item $n\geqq3$ について，$a_n$を漸化式で表せ．（解答経過の要点も，解答欄に記入せよ．）
            \end{myenud}
    \end{myenuc}


\end{document}
